\documentclass[12pt,a4paper]{article}
\usepackage[utf8]{inputenc}
\usepackage[portuguese]{babel}
\usepackage[T1]{fontenc}
\usepackage{geometry}
\usepackage{listings}
\usepackage{xcolor}
\usepackage{graphicx}

\geometry{margin=2.5cm}

\lstset{
    basicstyle=\ttfamily\small,
    keywordstyle=\color{blue},
    breaklines=true,
    frame=single,
    backgroundcolor=\color{gray!5}
}

\title{Manual de Instruções: Suíte de Segurança Informática}
\author{Paulo Abade - 23919}
\date{\today}

\begin{document}

\maketitle

\section{Pré-requisitos}
Para a correta execução desta suíte, o ambiente deve garantir os seguintes componentes:
\begin{itemize}
    \item \textbf{Sistema Operativo Recomendado:} Kali Linux
    \item \textbf{Python:} Versão 3.10 ou superior.
    \item \textbf{Dependências de Sistema:} Binário do \texttt{nmap}, \texttt{rsyslog} e \texttt{apache2} ativos no servidor.
\end{itemize}

\section{Instalação}
\subsection{Configuração Automática (Linux)}
O projeto inclui um script de automação para instalação de todas as bibliotecas Python necessárias (\texttt{cryptography}, \texttt{scapy}, \texttt{pyotp}, etc.).
\begin{lstlisting}[language=bash]
chmod +x install_deps.sh
./install_deps.sh
\end{lstlisting}


\section{Configuração do Servidor de Chat}
\begin{enumerate}
    \item No ficheiro \texttt{python-client.py}, altere a variável \texttt{SERVER\_IP} para o IP do seu servidor.
    \item Garanta que a porta \texttt{9999} está aberta na firewall do servidor:
    \begin{lstlisting}[language=bash]
    sudo ufw allow 9999/tcp
    \end{lstlisting}
\end{enumerate}

\section{Guia de Utilização}
O sistema é gerido através de um menu centralizado. Execute:
\begin{lstlisting}[language=bash]
python menu_cli.py
\end{lstlisting}

\subsection{Funcionalidades Principais}
\begin{itemize}
    \item \textbf{Opção 1 (Port Scan):} Insira o IP alvo e o intervalo de portas para auditoria.
    \item \textbf{Opção 2 (DoS Attack):} Forneça o IP alvo e o sistema sobrecarregará todas as portas TCP com pacotes.
    \item \textbf{Opção 3 (SYN Flood):} Forneça o IP alvo e a porta para um ataque de inundação SYN. 
    \item \textbf{Opção 4 (Análise de Logs):} Requer que a base de dados \texttt{GeoLite2-City.mmdb} esteja na pasta \texttt{CheckingFiles}.
    \item \textbf{Opção 5 (Servidor de Chat):} Deve ser iniciada no servidor antes de qualquer tentativa de ligação do cliente.
    \item \textbf{Opção 6 (Port Knocking):} Configure a sequência de portas e o IP do servidor para ativação.
    \item \textbf{Opção 7 (Gestor de Passwords):} No primeiro acesso, será gerado um ficheiro \texttt{2fa\_qr.png}. Leia este código com o Microsoft Authenticator 
    e/ou outra aplicação equivalente.
\end{itemize}

\section{Resolução de Problemas}
\begin{itemize}
    \item \textbf{Erro de Permissão (Scapy):} As opções de SYN Flood e Port Knocking requerem privilégios de administrador (\texttt{sudo python menu\_cli.py}).
    \item \textbf{Erro de Ficheiro:} Utilize a configuração inicial para garantir que os caminhos dos ficheiros estão corretos e que as dependências 
    estão instaladas.
\end{itemize}

\end{document}